% ---------------------------------------------------------------------------
%  fig1.tex -- Figure 1 of the CMOS chain paper, pure vector.
%
%  Replaces NewFigures/Figure1CMOS_v8.pdf (a PowerPoint export: 22 raster
%  objects, the density panel a 370x370 indexed bitmap at 218 ppi, and a
%  792x612 MediaBox hiding a superseded draft outside the CropBox).
%
%  Authored at exactly 510pt = \textwidth, so \includegraphics does no
%  rescaling: [width=\textwidth] in a figure*.  pdfinfo reads 508.09 bp.
%
%  Geometry note.  Every coordinate below is quoted in the ORIGINAL figure's
%  own units -- the 610.374 x 153.033 pt CropBox of Figure1CMOS_v8.pdf, read
%  out of its content stream -- and the x/y unit vectors carry the single
%  factor 510/610.374 = 0.8355542 that maps that space onto this canvas.  So
%  the numbers here can be checked one by one against the file being replaced.
%  Font sizes and line widths are absolute pt and so are NOT scaled.
%
%  Build: make fig1   (needs fig1_heatmap.dat, fig1_marginal.dat from
%  prep_fig1.py).
% ---------------------------------------------------------------------------
\documentclass[border=0pt]{standalone}
% ---------------------------------------------------------------------------
%  figstyle.tex -- shared style for every figure in the CMOS chain paper.
%  \input this from each figN.tex. Do not restyle locally; change it here.
%
%  Design rule: every figure is authored at EXACTLY the width it occupies on
%  the page, so \includegraphics applies no scaling and a 9pt label in one
%  figure is a 9pt label in every other. The paper's geometry, measured from
%  revtex4-2 [reprint,aps,pre]:
%      \columnwidth = 246.0pt = 3.404in   -> \figcolwidth
%      \textwidth   = 510.0pt = 7.057in   -> \figfullwidth
%  Include single-column figures with [width=\columnwidth] and full-width
%  figures (figure*) with [width=\textwidth]. Nothing else.
%
%  Verifying the width: pdfinfo reports big points (72/in), not LaTeX pt
%  (72.27/in). A correct column figure reads 245.08 bp; a correct full-width
%  figure reads 508.09 bp. Those are the numbers to check, not 246 and 510.
%
%  Fonts: Computer Modern, i.e. the LaTeX default. Load no font package.
%  Class sizes for reference: body 10, \small 9, \footnotesize 8, \scriptsize 7.
% ---------------------------------------------------------------------------
\usepackage{tikz}
\usepackage{pgfplots}
\usepackage{amsmath}
\usepackage{braket}
\pgfplotsset{compat=1.17}
\usetikzlibrary{arrows.meta,positioning,calc,backgrounds,fit,decorations.pathreplacing}
\usepgfplotslibrary{colormaps}

% ---- the two canonical widths -------------------------------------------
\newlength{\figcolwidth}   \setlength{\figcolwidth}{246.0pt}
\newlength{\figfullwidth}  \setlength{\figfullwidth}{510.0pt}

% ---- force an exact output width ----------------------------------------
%  Put \figcanvas{<width>}{<height>} as the FIRST line inside a tikzpicture.
%  It pins the bounding box, so the emitted PDF is exactly that size no
%  matter what the content does, and no rescaling happens at include time.
\newcommand{\figcanvas}[2]{\useasboundingbox (0,0) rectangle (#1,#2);}

% ---- type sizes (true points, because scale is 1) -----------------------
\newcommand{\figaxislabel}{\small}       % 9pt  -- axis labels
\newcommand{\figticklabel}{\footnotesize}% 8pt  -- tick labels
\newcommand{\figannot}{\footnotesize}    % 8pt  -- in-figure annotation
\newcommand{\figpanel}{\small\bfseries}  % 9pt  -- panel letters (a), (b)

% ---- palette -------------------------------------------------------------
\definecolor{figblueA}{HTML}{1F4E9C}  % dark blue   -- emphasis curve 1
\definecolor{figblueB}{HTML}{3C8DBC}  % mid blue    -- emphasis curve 2
\definecolor{figgreen}{HTML}{9CC7A9}  % pale green  -- de-emphasised curves
\definecolor{figgray}{HTML}{7F7F7F}   % neutral rules and guides
\definecolor{figred}{HTML}{C0392B}    % the one warning/annotation colour
\definecolor{fignodeA}{HTML}{6E8B74}  % node dot, k=0 side
\definecolor{fignodeB}{HTML}{1F4E9C}  % node dot, k=1 side
\definecolor{figshadeA}{HTML}{EAEFEA} % k=0 panel background
\definecolor{figshadeB}{HTML}{DCE8F5} % k=1 panel background
% second operator colour: the reduced MPO (physical dimension d), as against
% the full MPO (physical dimension D) drawn in figblueA/figblueB.
\definecolor{figtealA}{HTML}{1B6E6A} % reduced-MPO outline
\definecolor{figtealB}{HTML}{9FCFCB} % reduced-MPO fill

% ---- the V_dd ramp -------------------------------------------------------
%  Four samples of viridis at 0.08, 0.36, 0.64 and 0.92, the ramp the
%  matplotlib composite used, so a given V_dd is the same colour in every
%  panel that shows one.  The values are read out of the content stream of
%  PostReviewFigures/CompositeFigure3_v3.pdf, so they are matplotlib's own.
\definecolor{figvddA}{rgb}{0.283197,0.115680,0.436115}  % V_dd/V_T = 1.1
\definecolor{figvddB}{rgb}{0.179019,0.433756,0.557430}  % V_dd/V_T = 1.2
\definecolor{figvddC}{rgb}{0.170948,0.694384,0.493803}  % V_dd/V_T = 1.3
\definecolor{figvddD}{rgb}{0.793760,0.880678,0.120005}  % V_dd/V_T = 1.4

% ---- line weights --------------------------------------------------------
\pgfplotsset{
  figframe/.style={line width=0.6pt},
  figcurve/.style={line width=1.0pt},
  figemph/.style={line width=1.4pt},
}
\tikzset{
  figpanellabel/.style={font=\small\bfseries, inner sep=1pt},
  figwire/.style={line width=0.9pt},
  figlead/.style={figgray, line width=0.5pt, densely dotted},
  figarrow/.style={-{Stealth[length=4pt,width=3pt]}, line width=0.7pt},
}

% ---- the common axis style ----------------------------------------------
%  Matches the tick conventions already used by the matplotlib figures:
%  inward ticks on all four sides, minor ticks visible, no grid.
\pgfplotsset{
  figaxis/.style={
    scale only axis,
    axis line style={line width=0.6pt},
    tick align=inside,
    tickpos=both,
    minor tick num=1,
    tick label style={font=\figticklabel},
    label style={font=\figaxislabel},
    legend style={font=\figannot, draw=none, fill=none},
    every axis plot/.append style={figcurve},
  },
  figdensity/.style={
    figaxis,
    % `viridis high res' is matplotlib's own 256-entry table (the plain
    % `viridis' key is an 18-point approximation to it), so a density panel
    % ported from matplotlib keeps the colours it had.
    colormap/viridis high res,
    point meta min=0,
  },
}


% one x/y unit = one unit of the original figure's coordinate space
\newcommand{\figone}{0.8355542pt}

% ---- the MOSFET symbol ---------------------------------------------------
%  #1 x of the three-segment channel bar, #2 x of the gate plate (the two may
%  be in either order, which is how the pair gets mirrored), #3..#5 y of the
%  three channel stubs (far, bulk, near), #6 x the stubs run to, #7 y of the
%  gate lead, #8 x the gate lead runs to, #9 bulk-arrow tip x.
\newcommand{\figplate}[3]{\draw[figwire, line width=1.8pt] (#1,#2) -- (#1,#3);}
\begin{document}
\begin{tikzpicture}[x=\figone, y=\figone, every node/.style={inner sep=0pt}]
\figcanvas{\figfullwidth}{128pt}

% =========================================================================
%  1.  left panel -- the k=0 NOT gate
% =========================================================================
\fill[figshadeA] (3.31,3.14) rectangle (129.00,150.50);

% supply and node rails, all on x = 56.6 as three disjoint segments
\draw[figarrow] (56.6,103.43) -- (56.6,137.98);              % +V_dd
\draw[figarrow] (56.6,52.63)  -- (56.6,18.06);               % -V_dd
\draw[figwire]  (56.6,58.97)  -- (56.6,97.72);               % node rail
\node[anchor=base, font=\figaxislabel] at (58.0,141.6) {$V_{\rm dd}$};
\node[anchor=base, font=\figaxislabel] at (56.0,8.6)   {$-V_{\rm dd}$};

% PMOS (above) and NMOS (below): channel bar at 71.91, gate plate at 76.01
\draw[figwire] (71.91,93.18) -- (71.91,114.76);
\figplate{76.01}{93.18}{114.68}
\draw[figwire] (56.6,112.24) -- (71.91,112.24);
\draw[figwire] (56.6,104.41) -- (71.91,104.41);
\draw[figwire,-{Triangle[length=4pt,width=4pt]}] (71.91,104.41) -- (64.0,104.41);
\draw[figwire] (56.6,96.65)  -- (71.91,96.65);
\draw[figwire] (76.01,113.44) -- (99.02,113.44);

\draw[figwire] (71.91,41.40) -- (71.91,62.98);
\figplate{76.01}{41.40}{62.91}
\draw[figwire] (56.6,60.10) -- (71.91,60.10);
\draw[figwire] (56.6,51.69) -- (71.91,51.69);
\draw[figwire,-{Triangle[length=4pt,width=4pt]}] (63.9,51.69) -- (71.0,51.69);
\draw[figwire] (56.6,44.05) -- (71.91,44.05);
\draw[figwire] (76.01,42.31) -- (99.02,42.31);

\node[anchor=base, font=\figannot] at (85.5,119.0) {PMOS};
\node[anchor=base, font=\figannot] at (85.5,30.5)  {NMOS};

% output rail and the two nodes
\draw[figwire] (99.02,42.31) -- (99.02,113.44);
\draw[figwire] (36.21,75.5) -- (56.6,75.5);
\draw[figwire] (99.02,75.5) -- (117.82,75.5);
\fill[fignodeA] (36.21,75.5)  circle[radius=3.08pt];
\fill[fignodeB] (117.82,75.5) circle[radius=3.08pt];
\node[fignodeA, anchor=base, font=\figaxislabel] at (36.21,63.2)  {$v_0$};
\node[fignodeB, anchor=base, font=\figaxislabel] at (117.82,63.2) {$v_1$};

% rate arrows: up on the left of each pair, down on the right
\draw[figarrow] (27.7,101.0) -- (27.7,132.5);
\draw[figarrow] (32.0,132.5) -- (32.0,101.0);
\draw[figarrow] (27.7,21.0)  -- (27.7,52.5);
\draw[figarrow] (32.0,52.5)  -- (32.0,21.0);
\node[anchor=base east, font=\figaxislabel] at (26.0,109.6)  {$\lambda^{{\rm d},0}_{+}$};
\node[anchor=base west, font=\figaxislabel] at (33.9,109.6)  {$\lambda^{{\rm u},0}_{+}$};
\node[anchor=base east, font=\figaxislabel] at (26.0,31.6)   {$\lambda^{{\rm d},0}_{-}$};
\node[anchor=base west, font=\figaxislabel] at (33.9,31.6)   {$\lambda^{{\rm u},0}_{-}$};

% =========================================================================
%  2.  middle panel -- the CMOS unit as two logic gates
% =========================================================================
\fill[figshadeB] (162.97,82.05) rectangle (208.83,126.94);
\fill[figshadeA] (177.18,24.52) rectangle (223.04,69.41);
\draw[figwire] (159.18,45.43) rectangle (226.77,105.81);
\fill[white]   (168.56,88.54) -- (168.56,123.08) -- (205.39,105.81) -- cycle;
\draw[figwire] (168.56,88.54) -- (168.56,123.08) -- (205.39,105.81) -- cycle;
\fill[white]   (181.27,28.16) -- (181.27,62.69) -- (218.10,45.43) -- cycle;
\draw[figwire] (181.27,28.16) -- (181.27,62.69) -- (218.10,45.43) -- cycle;
% straddle each triangle's centre line, where it is widest
\node[figblueA, anchor=base west, font=\figannot] at (170.00,102.5) {$k=1$};
\node[figgray,  anchor=base west, font=\figannot] at (183.00,42.1)  {$k=0$};
\draw[figwire] (143.65,75.5) -- (159.18,75.5);
\draw[figwire] (226.77,75.5) -- (242.44,75.5);
\fill[fignodeA] (143.65,75.5) circle[radius=3.08pt];
\fill[fignodeB] (242.44,75.5) circle[radius=3.08pt];
\node[fignodeA, anchor=base, font=\figaxislabel] at (143.65,63.2) {$v_0$};
\node[fignodeB, anchor=base, font=\figaxislabel] at (242.44,63.2) {$v_1$};

% leaders: left panel <-> the k=0 gate, the k=1 gate <-> right panel
\draw[figlead] (129.00,4.07)   -- (177.97,25.19);
\draw[figlead] (208.68,126.94) -- (255.95,148.64);

% =========================================================================
%  3.  right panel -- the k=1 NOT gate, transistor pair mirrored
% =========================================================================
\fill[figshadeB] (257.96,3.14) rectangle (382.43,149.19);

\draw[figarrow] (328.88,103.43) -- (328.88,137.61);          % +V_dd
\draw[figarrow] (328.88,52.63)  -- (328.88,17.12);           % -V_dd
\draw[figwire]  (328.88,57.65)  -- (328.88,96.39);           % node rail
\node[anchor=base, font=\figaxislabel] at (330.3,141.0) {$V_{\rm dd}$};
\node[anchor=base, font=\figaxislabel] at (328.3,7.6)   {$-V_{\rm dd}$};

% mirrored: gate plate at 309.07, channel bar at 313.19
\figplate{309.07}{92.31}{113.81}
\draw[figwire] (313.19,92.31) -- (313.19,113.89);
\draw[figwire] (313.19,111.34) -- (328.88,111.34);
\draw[figwire] (313.19,103.43) -- (328.88,103.43);
\draw[figwire,-{Triangle[length=4pt,width=4pt]}] (313.19,103.43) -- (321.1,103.43);
\draw[figwire] (313.19,95.33)  -- (328.88,95.33);
\draw[figwire] (286.00,112.83) -- (309.07,112.83);

\figplate{309.07}{40.71}{62.22}
\draw[figwire] (313.19,40.38) -- (313.19,61.97);
\draw[figwire] (313.19,58.60) -- (328.88,58.60);
\draw[figwire] (313.19,50.67) -- (328.88,50.67);
\draw[figwire,-{Triangle[length=4pt,width=4pt]}] (321.1,50.67) -- (314.0,50.67);
\draw[figwire] (313.19,42.73) -- (328.88,42.73);
\draw[figwire] (286.00,41.70) -- (309.07,41.70);

\node[anchor=base, font=\figannot] at (307.0,119.0) {PMOS};
\node[anchor=base, font=\figannot] at (307.0,30.5)  {NMOS};

\draw[figwire] (286.00,41.70) -- (286.00,112.83);
\draw[figwire] (267.69,75.5) -- (286.00,75.5);
\draw[figwire] (328.88,75.5) -- (348.96,75.5);
\fill[fignodeA] (267.69,75.5) circle[radius=3.08pt];
\fill[fignodeB] (348.96,75.5) circle[radius=3.08pt];
\node[fignodeA, anchor=base, font=\figaxislabel] at (267.69,63.2) {$v_0$};
\node[fignodeB, anchor=base, font=\figaxislabel] at (348.96,63.2) {$v_1$};

\draw[figarrow] (355.1,101.0) -- (355.1,132.5);
\draw[figarrow] (359.4,132.5) -- (359.4,101.0);
\draw[figarrow] (355.1,21.0)  -- (355.1,52.5);
\draw[figarrow] (359.4,52.5)  -- (359.4,21.0);
\node[anchor=base east, font=\figaxislabel] at (353.4,109.6) {$\lambda^{{\rm d},1}_{+}$};
\node[anchor=base west, font=\figaxislabel] at (361.3,109.6) {$\lambda^{{\rm u},1}_{+}$};
\node[anchor=base east, font=\figaxislabel] at (353.4,31.6)  {$\lambda^{{\rm d},1}_{-}$};
\node[anchor=base west, font=\figaxislabel] at (361.3,31.6)  {$\lambda^{{\rm u},1}_{-}$};

% =========================================================================
%  4.  data panel -- single-unit steady state, and its anti-diagonal cut
% =========================================================================
%  65 x 65 vector cells: matrix plot* emits one filled rectangle per cell,
%  centred on its coordinate, so the panel spans -3.25 .. 3.25 (65 cells of
%  0.1) exactly as the original's pcolormesh did.  No bitmap anywhere.
% pgfplots' anchor arithmetic misbehaves under non-default x/y unit vectors,
% so the two data axes go in a default-unit scope and carry absolute pt
% positions; the original-figure units they came from are in the comments.
\draw[figgray, line width=1.4pt] (560.36,22.75) -- (560.36,144.04);
\node[figgray, font=\figannot, rotate=-90] at (552.7,123.0) {$v_0=-v_1$};
\node[anchor=base,      font=\figaxislabel] at (484.0,4.3)  {$v_1$};  % xlabel
\node[anchor=east,      font=\figaxislabel] at (397.0,84.5) {$v_0$};  % ylabel

\begin{scope}[x=1cm, y=1cm]
\begin{axis}[figdensity,
  at={(353.399pt,18.215pt)}, anchor=south west,   % orig (422.95, 21.80)
  axis on top,                              % 4225 filled cells would otherwise
                                            % bury the frame and the inward ticks
  width=102.00pt, height=102.94pt,                % orig 122.07 x 123.19
  xmin=-3.25, xmax=3.25, ymin=-3.25, ymax=3.25,
  xtick={-3.2,-1.6,0,1.6,3.2},              % nine tick marks as in the original,
  ytick={-3.2,-1.6,0,1.6,3.2},              % but only five carry a label: at the
  minor xtick={-2.4,-0.8,0.8,2.4},          % mandated 8pt, nine \"-3.2\"s would
  minor ytick={-2.4,-0.8,0.8,2.4},          % need 188pt of a 102pt axis
  xticklabel style={/pgf/number format/fixed, /pgf/number format/precision=1},
  yticklabel style={/pgf/number format/fixed, /pgf/number format/precision=1},
  % Ticks point inward, so without an explicit shift the labels sit against
  % the frame and the two "-3.2"s collide in the bottom-left corner.  The
  % canvas leaves only 2.9pt below the xlabel, hence the asymmetry: 1.5pt
  % down, 2.5pt out to the left where the gap to the schematic is wide.
  xticklabel shift=1.5pt,
  yticklabel shift=2.5pt,
  point meta max=0.006747483411018358,
]
\addplot[matrix plot*, mesh/cols=65, point meta=explicit]
   table[x=x, y=y, meta=z] {fig1_heatmap.dat};
\addplot[white, figcurve, forget plot] coordinates {(-3.25,3.25) (3.25,-3.25)};
\node[white, font=\figannot, rotate=-45, anchor=north west]
   at (axis cs:1.35,-1.35) {$v_0=-v_1$};
\end{axis}

\begin{axis}[
  at={(468.166pt,18.215pt)}, anchor=south west, scale only axis, hide axis,
  width=36pt, height=102.94pt,                    % orig zero line at 560.36
  xmin=0, xmax=1, ymin=-3.25, ymax=3.25,
]
\addplot[figblueA, figemph] table[x=p, y=v0] {fig1_marginal.dat};
\end{axis}
\end{scope}

\end{tikzpicture}
\end{document}
